\begin{exercise}{135}
Let $\mathcal C$ be an $\infty$-category and $x$ and object of $\mathcal C$. Show that the object
represented by $\Delta^0\to\mathcal C_{/x}$ which is adjoint to the map $\Delta^0\star\Delta^0\to\mathcal C$ given by $\mathrm{id}_x$ is a
terminal object.

\begin{proof}
Let $t$ be the specified object of $\mathcal C_{/x}$.
Associativity of joins gives
\[
 (\mathcal C_{/x})_{/t}\cong\mathcal C_{/\mathrm{id}_x}.
\]
Under this isomorphism, the slice projection is restriction to the final
vertex $1$ of the last $\Delta^1$; denote it by $v$.
Let $u\colon\mathcal C_{/\mathrm{id}_x}\to\mathcal C_{/x}$ instead be
restriction to vertex $0$.
Since $\mathrm{id}_x$ is degenerate, precomposition with
$\Delta^1\to\Delta^0$ gives a common section $s$ of $u$ and $v$:
$us=vs=\mathrm{id}$.

\setlength{\emergencystretch}{2em}
By Corollary~1.4.24(2), $u$ is a trivial fibration, hence a Joyal
equivalence by Proposition~2.2.12.
Two-out-of-three therefore makes $s$ and $v$ Joyal equivalences.
Corollary~1.4.24(1) makes $v$ a right fibration, so it is an
isofibration by Proposition~2.1.7.
Lemma~2.2.17 now makes $v$ a trivial fibration.
Lemma~4.1.2 proves that $t$ is terminal.
\end{proof}
\end{exercise}
